% Created 2024-11-08 vie 10:54
% Intended LaTeX compiler: pdflatex
\documentclass[11pt]{article}
\usepackage[utf8]{inputenc}
\usepackage[T1]{fontenc}
\usepackage{graphicx}
\usepackage{longtable}
\usepackage{wrapfig}
\usepackage{rotating}
\usepackage[normalem]{ulem}
\usepackage{amsmath}
\usepackage{amssymb}
\usepackage{capt-of}
\usepackage{hyperref}
\usepackage{minted}

\renewcommand{\contentsname}{Tabla de contenidos}
\usepackage{parskip}
\usepackage[utf8]{inputenc} %% For unicode chars
\usepackage{placeins}
\usepackage[margin=2.50cm]{geometry}
\usepackage[T1]{fontenc}
\usepackage{mathpazo}
\linespread{1.05}
\usepackage[scaled]{helvet}
\usepackage{courier}
\hypersetup{colorlinks=true,linkcolor=blue}
\RequirePackage{fancyvrb}
\DefineVerbatimEnvironment{verbatim}{Verbatim}{fontsize=\small,formatcom = {\color[rgb]{0.5,0,0}}}
\usepackage{mdframed}
\BeforeBeginEnvironment{minted}{\begin{mdframed}}
\AfterEndEnvironment{minted}{\end{mdframed}}
\setcounter{secnumdepth}{3}
\date{}
\title{Gestión de Biblioteca}
\hypersetup{
 pdfauthor={},
 pdftitle={Gestión de Biblioteca},
 pdfkeywords={},
 pdfsubject={},
 pdfcreator={Emacs 27.1 (Org mode 9.6.18)}, 
 pdflang={Spanish}}
\begin{document}

\maketitle
\setcounter{tocdepth}{3}
\tableofcontents

\setlength\parindent{10pt}
\section{Vamos a desarrollar el proyecto paso a paso}
\label{sec:org4aa0139}

Vamos a refactorizar el proyecto para usar servicios en lugar
de manejar toda la lógica en los controladores. A medida que avanzamos,
también terminaremos el \textbf{controlador de préstamos} y la lógica para
\textbf{devoluciones tardías}.

Voy a dividir el proceso en pasos claros:

\begin{enumerate}
\item \textbf{Crear los servicios}: Definiremos los servicios para cada entidad
(\texttt{LibroService}, \texttt{AutorService}, \texttt{UsuarioService},
\texttt{PrestamoService}).
\item \textbf{Refactorizar los controladores}: Actualizaremos cada controlador
para que delegue la lógica a los servicios.
\item \textbf{Implementar la lógica de devoluciones tardías}: Añadiremos esta
lógica en \texttt{PrestamoService} y actualizaremos \texttt{PrestamoController}
para manejar devoluciones y calcular si una devolución es tardía.
\end{enumerate}

\subsubsection{Paso 1: Crear los Servicios}
\label{sec:org16e4d3e}
\begin{itemize}
\item 1.1. \textbf{LibroService}
\label{sec:org3e8b7e3}
Este servicio manejará la lógica de negocio de los libros, como crear,
actualizar, eliminar y buscar libros. También implementaremos métodos de
paginación y filtro.

\begin{minted}[]{java}
import com.biblioteca.gestion.entities.Libro;
import com.biblioteca.gestion.repositories.LibroRepository;
import org.springframework.beans.factory.annotation.Autowired;
import org.springframework.data.domain.Page;
import org.springframework.data.domain.Pageable;
import org.springframework.stereotype.Service;

import java.util.Optional;

@Service
public class LibroService {

    @Autowired
    private LibroRepository libroRepository;

    public Page<Libro> listarLibros(String titulo, String genero,
                                     Pageable pageable) {
        if (titulo != null) {
            return libroRepository.findByTituloContaining(titulo, pageable);
        } else if (genero != null) {
            return libroRepository.findByGenero(genero, pageable);
        } else {
            return libroRepository.findAll(pageable);
        }
    }

    public Optional<Libro> obtenerLibro(Long id) {
        return libroRepository.findById(id);
    }

    public Libro crearLibro(Libro libro) {
        return libroRepository.save(libro);
    }

    public Optional<Libro> actualizarLibro(Long id, Libro libroActualizado) {
        return libroRepository.findById(id).map(libro -> {
            libro.setTitulo(libroActualizado.getTitulo());
            libro.setGenero(libroActualizado.getGenero());
            libro.setAñoPublicacion(libroActualizado.getAñoPublicacion());
            libro.setEstado(libroActualizado.getEstado());
            libro.setAutor(libroActualizado.getAutor());
            return libroRepository.save(libro);
        });
    }

    public boolean eliminarLibro(Long id) {
        if (libroRepository.existsById(id)) {
            libroRepository.deleteById(id);
            return true;
        }
        return false;
    }
}
\end{minted}

\item 1.2. \textbf{AutorService}
\label{sec:org2d2f949}
El servicio de autores se encarga de manejar las operaciones CRUD para
la entidad \texttt{Autor}.

\begin{minted}[]{java}
import com.biblioteca.gestion.entities.Autor;
import com.biblioteca.gestion.repositories.AutorRepository;
import org.springframework.beans.factory.annotation.Autowired;
import org.springframework.stereotype.Service;

import java.util.List;
import java.util.Optional;

@Service
public class AutorService {

    @Autowired
    private AutorRepository autorRepository;

    public List<Autor> listarAutores() {
        return autorRepository.findAll();
    }

    public Autor crearAutor(Autor autor) {
        return autorRepository.save(autor);
    }

    public Optional<Autor> actualizarAutor(Long id, Autor autorActualizado) {
        return autorRepository.findById(id).map(autor -> {
            autor.setNombre(autorActualizado.getNombre());
            autor.setNacionalidad(autorActualizado.getNacionalidad());
            return autorRepository.save(autor);
        });
    }

    public boolean eliminarAutor(Long id) {
        if (autorRepository.existsById(id)) {
            autorRepository.deleteById(id);
            return true;
        }
        return false;
    }
}
\end{minted}

\item 1.3. \textbf{UsuarioService}
\label{sec:org7acf145}
Este servicio manejará la lógica relacionada con los usuarios.

\begin{minted}[]{java}
import com.biblioteca.gestion.entities.Usuario;
import com.biblioteca.gestion.repositories.UsuarioRepository;
import org.springframework.beans.factory.annotation.Autowired;
import org.springframework.stereotype.Service;

import java.util.Optional;

@Service
public class UsuarioService {

    @Autowired
    private UsuarioRepository usuarioRepository;

    public Usuario registrarUsuario(Usuario usuario) {
        return usuarioRepository.save(usuario);
    }

    public Optional<Usuario> obtenerUsuario(Long id) {
        return usuarioRepository.findById(id);
    }

    public Optional<Usuario> obtenerUsuarioPorEmail(String email) {
        return usuarioRepository.findByEmail(email);
    }
}
\end{minted}

\item 1.4. \textbf{PrestamoService} (Incluyendo la Lógica de Devoluciones Tardías)
\label{sec:org2055122}
Aquí es donde implementamos la lógica de los préstamos, incluyendo la
verificación de \textbf{devoluciones tardías}.

\begin{minted}[]{java}
import com.biblioteca.gestion.entities.Libro;
import com.biblioteca.gestion.entities.Prestamo;
import com.biblioteca.gestion.repositories.LibroRepository;
import com.biblioteca.gestion.repositories.PrestamoRepository;
import org.springframework.beans.factory.annotation.Autowired;
import org.springframework.stereotype.Service;

import java.time.LocalDate;
import java.util.List;
import java.util.Optional;

@Service
public class PrestamoService {

    @Autowired
    private PrestamoRepository prestamoRepository;

    @Autowired
    private LibroRepository libroRepository;

    public Prestamo registrarPrestamo(Prestamo prestamo) {
        Optional<Libro> libro = libroRepository.findById(prestamo.getLibro().getId());
        if (libro.isPresent() && "disponible".equalsIgnoreCase(libro.get().getEstado())) {
            libro.get().setEstado("prestado");
            prestamo.setFechaPrestamo(LocalDate.now());
            return prestamoRepository.save(prestamo);
        }
        throw new IllegalArgumentException("El libro no está disponible para préstamo.");
    }

    public Optional<Prestamo> devolverLibro(Long prestamoId) {
        return prestamoRepository.findById(prestamoId).map(prestamo -> {
            prestamo.setFechaDevolucion(LocalDate.now());
            Libro libro = prestamo.getLibro();
            libro.setEstado("disponible");
            // Verificar si la devolución es tardía
            if (prestamo.getFechaPrestamo().plusDays(14).isBefore(LocalDate.now())) {
                prestamo.setDevolucionTardia(true); // Marcar devolución tardía
            }
            prestamoRepository.save(prestamo);
            return prestamo;
        });
    }

    public List<Prestamo> listarPrestamosPorUsuario(Long usuarioId) {
        return prestamoRepository.findByUsuarioId(usuarioId);
    }
}
\end{minted}
\end{itemize}

\subsubsection{Paso 2: Refactorizar los Controladores}
\label{sec:org4ff3eff}
Ahora, actualizaremos los controladores para delegar la lógica en estos
servicios. Por ejemplo, \texttt{PrestamoController}:

\begin{minted}[]{java}
import com.biblioteca.gestion.entities.Prestamo;
import com.biblioteca.gestion.services.PrestamoService;
import org.springframework.beans.factory.annotation.Autowired;
import org.springframework.http.ResponseEntity;
import org.springframework.web.bind.annotation.*;

import java.util.List;

@RestController
@RequestMapping("/prestamos")
public class PrestamoController {

    @Autowired
    private PrestamoService prestamoService;

    @PostMapping
    public ResponseEntity<Prestamo> registrarPrestamo(@RequestBody Prestamo prestamo) {
        try {
            Prestamo nuevoPrestamo = prestamoService.registrarPrestamo(prestamo);
            return ResponseEntity.ok(nuevoPrestamo);
        } catch (IllegalArgumentException e) {
            return ResponseEntity.badRequest().body(null);
        }
    }

    @PutMapping("/{id}/devolver")
    public ResponseEntity<Prestamo> devolverLibro(@PathVariable Long id) {
        return prestamoService.devolverLibro(id)
                .map(ResponseEntity::ok)
                .orElse(ResponseEntity.notFound().build());
    }

    @GetMapping("/usuario/{usuarioId}")
    public ResponseEntity<List<Prestamo>> listarPrestamosPorUsuario(@PathVariable Long usuarioId) {
        List<Prestamo> prestamos = prestamoService.listarPrestamosPorUsuario(usuarioId);
        return ResponseEntity.ok(prestamos);
    }
}
\end{minted}

\subsubsection{Resumen}
\label{sec:org9d1db53}
\begin{enumerate}
\item \textbf{Servicios creados} para \texttt{Libro}, \texttt{Autor}, \texttt{Usuario}, y \texttt{Prestamo},
encapsulando la lógica de negocio.
\item \textbf{Controladores refactorizados} para delegar en los servicios.
\item \textbf{Lógica de devoluciones tardías} implementada en \texttt{PrestamoService} y
gestionada desde \texttt{PrestamoController}.
\end{enumerate}

Con estos cambios, la estructura es más clara y permite una gestión de
lógica de negocio eficiente y separada de la capa de controladores.
¡Listo para revisar cualquier duda o continuar!
\end{document}
